\documentclass[a4paper,12pt]{article}

\usepackage[T2A]{fontenc}
\usepackage[utf8]{inputenc}
\usepackage[russian]{babel}
\usepackage{amsmath,amssymb,mathtools}
\usepackage{PTSans}
\renewcommand{\familydefault}{\sfdefault}
\usepackage{icomma}
\usepackage[a4paper,margin=20mm,headheight=25pt]{geometry}
\usepackage{microtype}
\usepackage{tikz}
\usetikzlibrary{calc,arrows.meta,positioning,shapes.geometric}
\usepackage{tabularx,booktabs}
\usepackage{enumitem}
\usepackage{parskip}
\usepackage{fancyhdr}
\usepackage{setspace}
\usepackage{xcolor}

\definecolor{accent}{HTML}{008080}
\definecolor{darkaccent}{HTML}{004D4D}
\definecolor{textgray}{HTML}{333333}
\definecolor{softgray}{HTML}{F2F5F5}

\color{textgray}
\onehalfspacing
\setlength{\parindent}{0pt}
\setlist[itemize]{leftmargin=1.4em,itemsep=0.25em,topsep=0.35em}
\setlist[enumerate]{leftmargin=1.7em,itemsep=0.45em,topsep=0.4em}
\renewcommand{\arraystretch}{1.28}

\pagestyle{fancy}
\fancyhf{}
\fancyhead[L]{\small\textcolor{darkaccent}{Физика: равноускоренное движение}}
\fancyhead[R]{\small\textcolor{darkaccent}{Володя Хачатурян}}
\fancyfoot[C]{\thepage}
\renewcommand{\headrulewidth}{0.4pt}
\renewcommand{\headrule}{%
  \hbox to\headwidth{%
    \color{accent}\leaders\hrule height \headrulewidth\hfill
  }%
}

\newcommand{\ms}{\,\text{м/с}}
\newcommand{\mss}{\,\text{м/с}^2}
\newcommand{\important}[1]{\textcolor{darkaccent}{\textbf{#1}}}
\newcommand{\formula}[1]{\[\boxed{#1}\]}

\tikzset{
  startstop/.style={
    ellipse,
    draw=accent,
    line width=0.9pt,
    align=center,
    text width=7.2cm,
    minimum height=0.9cm,
    inner sep=1.5pt,
    fill=softgray
  },
  process/.style={
    rectangle,
    rounded corners=4pt,
    draw=accent,
    line width=0.9pt,
    align=center,
    text width=5.8cm,
    minimum height=1.0cm,
    inner sep=2.5pt,
    fill=white
  },
  decision/.style={
    diamond,
    aspect=1.8,
    draw=accent,
    line width=0.9pt,
    align=center,
    text width=3.8cm,
    inner sep=0.8pt,
    fill=softgray
  },
  arrow/.style={
    -{Stealth[length=3mm,width=2mm]},
    line width=0.9pt,
    color=darkaccent
  },
  branch/.style={
    font=\small\bfseries,
    text=darkaccent,
    fill=white,
    inner sep=1pt
  }
}

\begin{document}

% =========================================================
% ТИТУЛЬНЫЙ ЛИСТ
% =========================================================

\begin{titlepage}
\thispagestyle{empty}
\centering
\vspace*{2.0cm}

{\Large\textcolor{accent}{Чек-лист}\par}
\vspace{0.8cm}

{\huge\bfseries
Прямолинейное равноускоренное движение
\par}

\vspace{0.35cm}

{\Large
Путь, скорость, ускорение и чтение условия задачи
\par}

\vspace{1.4cm}

{\large ОГЭ по физике\par}

\vspace{0.5cm}

{\large \today\par}

\vfill

{\large\bfseries
Лёвин Артём Александрович
\par}

\vspace{0.15cm}

{\large
эксклюзивно для Володи Хачатуряна
\par}

\vspace{1.4cm}

\begin{tikzpicture}[remember picture,overlay]
  \draw[
    accent,
    line width=1.2pt
  ]
    ([xshift=18mm,yshift=16mm]current page.south west)
    ..
    controls
    ([xshift=58mm,yshift=32mm]current page.south west)
    and
    ([xshift=-58mm,yshift=5mm]current page.south east)
    ..
    ([xshift=-18mm,yshift=19mm]current page.south east);
\end{tikzpicture}

\end{titlepage}

% =========================================================
% ОСНОВНОЙ ЧЕК-ЛИСТ
% =========================================================

\section*{Основной чек-лист}

\subsection*{1. Что означает равноускоренное прямолинейное движение}

Равноускоренное прямолинейное движение --- движение по прямой,
при котором проекция ускорения постоянна:

\[
a_x=\text{const}.
\]

Если ось $Ox$ направлена по движению тела, то:

\begin{itemize}
  \item при разгоне $a_x>0$;
  \item при торможении $a_x<0$;
  \item фразы
  «трогается с места»,
  «начинает движение из состояния покоя»,
  «стартует из состояния покоя»
  означают
  \[
  v_{0x}=0;
  \]
  \item если такого указания в условии нет, автоматически считать
  $v_{0x}=0$ нельзя.
\end{itemize}

\subsection*{2. Четыре главные формулы}

Для движения вдоль одной оси при $a_x=\text{const}$:

\begin{align*}
v_x &= v_{0x}+a_xt,\\[2mm]
s &= v_{0x}t+\frac{a_xt^2}{2},\\[2mm]
s &= \frac{v_{0x}+v_x}{2}\,t,\\[2mm]
v_x^2-v_{0x}^2 &= 2a_xs.
\end{align*}

Здесь $s$ --- перемещение вдоль выбранной оси.
В рассматриваемых задачах тело движется по прямой в одном направлении
до момента остановки, поэтому численно перемещение совпадает
с пройденным путём.

\important{Как выбирать формулу.}
Сначала выпишите известные величины и то, что требуется найти.
После этого выберите формулу, в которой присутствует искомая величина
и как можно меньше других неизвестных.

\subsection*{3. Единицы СИ перед подстановкой}

Перед вычислениями удобно сразу перейти к СИ:

\[
1\,\text{км}=1000\,\text{м},
\qquad
1\,\text{см}=0{,}01\,\text{м}.
\]

Например,

\[
15\,\text{см}=0{,}15\,\text{м},
\qquad
5\,\text{км}=5000\,\text{м}.
\]

\important{Проверка.}
Если в одной формуле одновременно встречаются метры и сантиметры
или метры и километры, единицы ещё требуется привести к общей системе.

\subsection*{4. Движение из состояния покоя}

Если

\[
v_{0x}=0,
\]

формула пути упрощается:

\formula{s=\frac{a_xt^2}{2}}

\textbf{Пример 1.}
Девочка начинает движение из состояния покоя
с ускорением $1{,}6\mss$.
Какой путь она пройдёт за первые $4$ с?

\[
s=
\frac{1{,}6\cdot 4^2}{2}
=
\frac{1{,}6\cdot16}{2}
=
12{,}8\,\text{м}.
\]

\textbf{Ответ:}
$12{,}8$ м.

\medskip

\textbf{Пример 2.}
Через $10$ с после старта ракета находится на расстоянии
$5$ км от поверхности Земли.
Считать, что ракета начала движение из состояния покоя
и двигалась с постоянным ускорением.

\[
s=5000\,\text{м},
\qquad
t=10\,\text{с},
\qquad
v_{0x}=0.
\]

Используем

\[
s=\frac{a_xt^2}{2}.
\]

Тогда

\[
5000=\frac{a_x\cdot10^2}{2}.
\]

Отсюда

\[
a_x=
\frac{2\cdot5000}{100}
=
100\mss.
\]

\textbf{Ответ:}
$100\mss$.

\subsection*{5. Если требуется найти ускорение}

При $v_{0x}=0$ удобно сразу выразить ускорение:

\formula{a_x=\frac{2s}{t^2}}

Например, каретка проходит $15$ см за $0{,}26$ с,
начиная движение из состояния покоя.

Сначала:

\[
15\,\text{см}=0{,}15\,\text{м}.
\]

Затем:

\[
a_x=
\frac{2\cdot0{,}15}{0{,}26^2}
\approx
4{,}44\mss.
\]

\important{Полезный навык.}
Сначала выразите искомую величину из формулы,
затем подставляйте числа.
Так решение легче проверять.

\subsection*{6. Начальная скорость может быть неизвестной}

Пусть лыжник проходит $50$ м за $10$ с
с постоянным ускорением $0{,}2\mss$.

В условии отсутствует указание,
что движение началось из состояния покоя.
Поэтому $v_{0x}$ следует считать неизвестной величиной.

Используем:

\[
s=v_{0x}t+\frac{a_xt^2}{2}.
\]

Подставляем данные:

\[
50
=
10v_{0x}
+
\frac{0{,}2\cdot10^2}{2}.
\]

Получаем:

\[
50=10v_{0x}+10,
\]

\[
10v_{0x}=40,
\]

\[
v_{0x}=4\ms.
\]

Теперь найдём скорость в конце участка:

\[
v_x=v_{0x}+a_xt.
\]

\[
v_x
=
4+0{,}2\cdot10
=
6\ms.
\]

\textbf{Ответ:}
$6\ms$.

\important{Контроль смысла.}
Если выбранная формула вообще не содержит искомую величину,
возможно, она нужна для нахождения промежуточной величины.
В подобных задачах решение может состоять из двух этапов.

\subsection*{7. Торможение и знак ускорения}

Если ось $Ox$ направлена по движению тела,
а скорость уменьшается, то ускорение имеет отрицательную проекцию:

\[
a_x<0.
\]

Например, автомобиль движется со скоростью $15\ms$
и тормозит с ускорением

\[
a_x=-5\mss
\]

до полной остановки.

Удобно использовать формулу

\[
v_x^2-v_{0x}^2=2a_xs.
\]

В конце торможения

\[
v_x=0.
\]

Поэтому

\[
0^2-15^2
=
2\cdot(-5)s.
\]

\[
-225=-10s.
\]

\[
s=22{,}5\,\text{м}.
\]

\textbf{Ответ:}
$22{,}5$ м.

\important{Физический смысл знака.}
Минус у $a_x$ показывает,
что ускорение направлено противоположно скорости.

\subsection*{8. Один и тот же тип формулы для разных объектов}

Автомобиль, лыжник, каретка, ракета и пуля
могут описываться одними и теми же кинематическими формулами.

Например, если пуля разгоняется внутри ствола
из состояния покоя,
её конечная скорость равна $250\ms$,
а длина ствола равна $0{,}10$ м,
то используется та же связь:

\[
v_x^2-v_{0x}^2=2a_xs.
\]

Так как

\[
v_{0x}=0,
\]

получаем:

\[
250^2
=
2a_x\cdot0{,}10.
\]

Отсюда

\[
a_x=312500\mss.
\]

Главное при выборе формулы --- набор известных и неизвестных величин,
а не конкретный объект из текста задачи.

\subsection*{9. Сложная задача на тормозной путь: разбиваем процесс на этапы}

Пусть на последнем километре тормозного пути
скорость поезда уменьшилась на $10\ms$,
а в конце поезд полностью остановился.

На начале последнего километра скорость равна:

\[
v_{0x}=10\ms.
\]

Для последнего километра:

\[
s=1000\,\text{м},
\qquad
v_x=0.
\]

Используем:

\[
v_x^2-v_{0x}^2=2a_xs.
\]

Тогда

\[
0^2-10^2
=
2a_x\cdot1000.
\]

\[
-100=2000a_x.
\]

\[
a_x=-0{,}05\mss.
\]

Теперь пусть весь тормозной путь равен $4$ км:

\[
s_{\text{торм}}=4000\,\text{м}.
\]

Ускорение в течение всего равнозамедленного движения постоянно,
поэтому снова используем

\[
v_x^2-v_{0x}^2=2a_xs.
\]

Для всего торможения:

\[
0-v_{0x}^2
=
2\cdot(-0{,}05)\cdot4000.
\]

\[
-v_{0x}^2=-400.
\]

\[
v_{0x}^2=400.
\]

Скорость --- неотрицательная величина, поэтому

\[
v_{0x}=20\ms.
\]

\textbf{Ответ:}
$20\ms$.

\important{Идея.}
Сначала используйте участок,
на котором достаточно данных для нахождения ускорения.
После этого найденное постоянное ускорение можно использовать
для всего движения.

\subsection*{10. «Скорость уменьшилась на» и сама скорость}

Фраза

\[
\text{«скорость уменьшилась на }10\ms\text{»}
\]

описывает изменение скорости:

\[
|\Delta v|=10\ms.
\]

Это ещё не означает,
что начальная или конечная скорость обязательно равна $10\ms$.

В предыдущем примере получить скорость $10\ms$
на начале последнего километра удалось потому,
что в конце этого километра поезд полностью остановился:

\[
v_{\text{кон}}=0.
\]

Следовательно,

\[
v_{\text{нач}}-0=10\ms.
\]

\subsection*{11. «За $n$ секунд» и «за $n$-ю секунду»}

Это два разных типа вопросов.

Если $v_{0x}=0$, то путь от начала движения за время $t$:

\[
s(t)=\frac{a_xt^2}{2}.
\]

Поэтому пути,
пройденные \emph{от начала движения}
за $1$, $2$, $3$, $4$ секунды,
относятся как

\[
1^2:2^2:3^2:4^2
=
1:4:9:16.
\]

Например,

\[
\frac{s(3\,\text{с})}{s(1\,\text{с})}
=
\frac{\frac{a_x\cdot3^2}{2}}
{\frac{a_x\cdot1^2}{2}}
=
9.
\]

Значит, путь за первые $3$ секунды
в $9$ раз больше пути за первую секунду.

\subsection*{12. Путь только за отдельную секунду}

Если требуется путь именно за $n$-ю секунду,
нужно из пути за первые $n$ секунд
вычесть путь за первые $n-1$ секунд:

\formula{\Delta s_n=s(n)-s(n-1)}

Пусть

\[
\tau=1\,\text{с}
\]

--- длительность одной секунды.

При $v_{0x}=0$:

\begin{align*}
\Delta s_n
&=
\frac{a_x(n\tau)^2}{2}
-
\frac{a_x((n-1)\tau)^2}{2}
\\[1mm]
&=
\frac{a_x\tau^2}{2}
\left(
n^2-(n-1)^2
\right)
\\[1mm]
&=
\frac{a_x\tau^2}{2}
(2n-1).
\end{align*}

Для пятой секунды:

\[
\Delta s_5
=
s(5\,\text{с})
-
s(4\,\text{с}).
\]

При $v_{0x}=0$:

\[
\Delta s_5
=
\frac{a_x(1\,\text{с})^2}{2}
(25-16).
\]

Следовательно,

\[
\Delta s_5
=
4{,}5a_x(1\,\text{с})^2.
\]

Если

\[
a_x=2\mss,
\]

то путь только за десятую секунду:

\[
\Delta s_{10}
=
\frac{2\cdot1^2}{2}
(2\cdot10-1)
=
19\,\text{м}.
\]

\important{Полезное наблюдение.}
При $v_{0x}=0$ и постоянном ускорении
пути за последовательные отдельные секунды
относятся как

\[
1:3:5:7:9:\ldots
\]

Суммарные пути от начала движения
к моментам $1$, $2$, $3$, $4$, $\ldots$ с
относятся как

\[
1:4:9:16:\ldots
\]

% =========================================================
% БЛОК-СХЕМА
% =========================================================

\clearpage
\thispagestyle{empty}

\begin{center}

{\Large\bfseries\textcolor{darkaccent}{
Алгоритм: решение задачи на равноускоренное движение
}\par}

\vspace{2mm}

\begin{tikzpicture}[
  font=\small,
  node distance=4mm and 8mm
]

\node[startstop] (start)
{Прочитать условие и выбрать положительное направление оси $Ox$};

\node[process,below=of start] (data)
{Выписать данные и искомую величину.
Перевести длины и время в СИ};

\node[decision,below=8mm of data] (rest)
{Есть указание
«из состояния покоя»
или
«трогается с места»?};

\node[
  process,
  text width=4.2cm
] (zero)
at ($(rest.south)+(-36mm,-13mm)$)
{Записать $v_{0x}=0$};

\node[
  process,
  text width=4.2cm
] (keep)
at ($(rest.south)+(36mm,-13mm)$)
{Оставить $v_{0x}$ известной
или неизвестной величиной
по условию};

\node[decision] (brake)
at ($(rest.south)+(0,-42mm)$)
{Тело замедляется при оси,
направленной по движению?};

\node[
  process,
  text width=4.2cm
] (negative)
at ($(brake.south)+(-36mm,-13mm)$)
{Использовать $a_x<0$};

\node[
  process,
  text width=4.2cm
] (sign)
at ($(brake.south)+(36mm,-13mm)$)
{Определить знак $a_x$
по направлению ускорения};

\node[process] (choose)
at ($(brake.south)+(0,-42mm)$)
{Выбрать формулу,
содержащую искомую величину
и минимум неизвестных};

\node[process,below=of choose] (solve)
{Подставить данные,
решить уравнение,
проверить единицы
и физический смысл знака};

\node[startstop,below=of solve] (finish)
{Записать ответ с единицами измерения};

\draw[arrow] (start) -- (data);
\draw[arrow] (data) -- (rest);

\draw[arrow]
(rest)
--
node[branch,above left]{да}
(zero);

\draw[arrow]
(rest)
--
node[branch,above right]{нет}
(keep);

\draw[arrow]
(zero)
|-
(brake.west);

\draw[arrow]
(keep)
|-
(brake.east);

\draw[arrow]
(brake)
--
node[branch,above left]{да}
(negative);

\draw[arrow]
(brake)
--
node[branch,above right]{нет / разгон}
(sign);

\draw[arrow]
(negative)
|-
(choose.west);

\draw[arrow]
(sign)
|-
(choose.east);

\draw[arrow] (choose) -- (solve);
\draw[arrow] (solve) -- (finish);

\end{tikzpicture}

\end{center}

\clearpage

% =========================================================
% ТИПИЧНЫЕ ОШИБКИ
% =========================================================

\section*{Типичные ошибки и самопроверка}

\begin{enumerate}

  \item
  \textbf{Автоматически писать $v_{0x}=0$.}

  Делайте это только при прямом указании
  на начало движения из состояния покоя.

  \item
  \textbf{Забывать знак ускорения при торможении.}

  Если ось направлена по движению,
  то при замедлении

  \[
  a_x<0.
  \]

  \item
  \textbf{Подставлять сантиметры и километры
  вместе с метрами.}

  Сначала приведите величины к СИ.

  \item
  \textbf{Путать изменение скорости с самой скоростью.}

  Запись

  \[
  |\Delta v|=10\ms
  \]

  означает изменение скорости на $10\ms$.
  Значение одной из скоростей определяется
  дополнительной информацией из условия.

  \item
  \textbf{Путать «за $10$ секунд»
  и «за десятую секунду».}

  Путь за первые десять секунд:

  \[
  s(10\,\text{с}).
  \]

  Путь только за десятую секунду:

  \[
  s(10\,\text{с})-s(9\,\text{с}).
  \]

  \item
  \textbf{Использовать формулу,
  в которой слишком много неизвестных.}

  Сравните несколько формул
  и выберите наиболее информативную.

  \item
  \textbf{Терять промежуточный этап.}

  Иногда сначала приходится найти
  $v_{0x}$ или $a_x$,
  после чего переходить к искомой величине.

  \item
  \textbf{Забывать ответ.}

  После вычислений укажите
  численное значение
  и единицы измерения.

\end{enumerate}

% =========================================================
% ПАМЯТКА
% =========================================================

\section*{Короткая памятка по формулам}

\begin{center}

\begin{tabularx}{\linewidth}{
  @{}
  >{\bfseries}p{0.40\linewidth}
  X
  @{}
}

\toprule

Что связывает формула
&
Формула
\\

\midrule

Скорость, ускорение, время
&
$v_x=v_{0x}+a_xt$
\\

Путь, начальная скорость, ускорение, время
&
$s=v_{0x}t+\dfrac{a_xt^2}{2}$
\\

Путь, начальная и конечная скорости, время
&
$s=\dfrac{v_{0x}+v_x}{2}\,t$
\\

Путь, скорости, ускорение
&
$v_x^2-v_{0x}^2=2a_xs$
\\

Путь только за $n$-ю секунду
&
$\Delta s_n=s(n)-s(n-1)$
\\

\bottomrule

\end{tabularx}

\end{center}

% =========================================================
% ТРЕНИРОВОЧНЫЙ БЛОК
% =========================================================

\clearpage

\section*{Тренировочный блок --- Путь А}

Решите задачи письменно.

Для каждой задачи:

\begin{itemize}
  \item выпишите «Дано»;
  \item выберите положительное направление оси;
  \item определите знаки проекций;
  \item переведите величины в СИ;
  \item выберите формулу;
  \item выполните вычисления;
  \item запишите ответ с единицами измерения.
\end{itemize}

\begin{enumerate}

  \item
  Тело начинает движение из состояния покоя
  с постоянным ускорением

  \[
  1{,}4\mss.
  \]

  Какой путь оно пройдёт за $6$ с?

  \vspace{8mm}

  \item
  Каретка проходит $32$ см за $0{,}40$ с,
  начиная движение из состояния покоя.

  Найдите модуль ускорения каретки.

  \vspace{8mm}

  \item
  Лыжник проходит по прямому участку $72$ м
  за $8$ с с постоянным ускорением

  \[
  1\mss.
  \]

  Начальная скорость в условии не задана.

  Найдите скорость лыжника в конце участка.

  \vspace{8mm}

  \item
  Тело начинает движение из состояния покоя
  с постоянным ускорением

  \[
  1{,}6\mss.
  \]

  Какой путь оно проходит
  только за седьмую секунду движения?

  \vspace{8mm}

  \item
  Поезд тормозит равнозамедленно
  до полной остановки.

  Весь тормозной путь равен $2{,}5$ км.

  На последних $900$ м
  его скорость уменьшилась на $12\ms$.

  Найдите:

  \begin{enumerate}[label=\alph*)]
    \item ускорение поезда;
    \item скорость поезда в начале торможения.
  \end{enumerate}

\end{enumerate}

% =========================================================
% ОТВЕТЫ
% =========================================================

\clearpage

\section*{Ответы}

\begin{center}

\begin{tabularx}{0.88\linewidth}{
  @{}
  c
  X
  @{}
}

\toprule

\textbf{№}
&
\textbf{Ответ}
\\

\midrule

1
&
$25{,}2\,\text{м}$
\\

2
&
$4\mss$
\\

3
&
$13\ms$
\\

4
&
$10{,}4\,\text{м}$
\\

5
&
$a_x=-0{,}08\mss$,
\qquad
$v_{0x}=20\ms$
\\

\bottomrule

\end{tabularx}

\end{center}

\vfill

\begin{center}
\small
\textcolor{darkaccent}{
Ключевой принцип:
сначала понять физический смысл условия,
затем выбирать формулу.
}
\end{center}

\end{document}